\usetheme{default}
\usecolortheme{default}
\usepackage{xcolor}
\definecolor{courseaccent}{HTML}{A02B0C}
\definecolor{courserule}{HTML}{C4BCAB}
\setbeamercolor{structure}{fg=courseaccent}
\setbeamercolor{frametitle}{fg=courseaccent}
\setbeamercolor{framesubtitle}{fg=gray}
\setbeamerfont{frametitle}{size=\Large,series=\bfseries}
\setbeamerfont{framesubtitle}{size=\small}
\setbeamertemplate{navigation symbols}{}
\setbeamertemplate{footline}{%
  \hfill\insertframenumber/\inserttotalframenumber\hspace{2em}\vspace{1ex}}
\setbeamertemplate{itemize items}[circle]
\renewcommand{\arraystretch}{1.2}
\AtBeginEnvironment{longtable}{\small}
\newcommand{\codesize}{\footnotesize}
\fvset{fontsize=\codesize}
% 図は高さを抑えて、同じフレームの文章と一緒に収める。キャプションの「Figure 1:」は付けない
\makeatletter
\def\maxwidth{\ifdim\Gin@nat@width>0.8\linewidth 0.8\linewidth\else\Gin@nat@width\fi}
\def\maxheight{\ifdim\Gin@nat@height>0.30\textheight 0.30\textheight\else\Gin@nat@height\fi}
\makeatother
\setkeys{Gin}{width=\maxwidth,height=\maxheight,keepaspectratio}
\setbeamertemplate{caption}{\centering\scriptsize\insertcaption\par}
\setlength{\parskip}{0pt}
% 長いコードの行は、はみ出さずに折り返す
\usepackage{fvextra}
\DefineVerbatimEnvironment{Highlighting}{Verbatim}{breaklines,commandchars=\\\{\}}
